% =====================================================================
%  26.3  Applications -- the two workhorse families.
%  Exports: sec:c30-apps  thm:c30-pairsum  lst:c30-pairsum
%  WORKED examples only -- NO solvedproblem (in-section solved stays 2).
%  Both families run on the running example T and S={4,8,6,10} and are
%  hand-checked against direct pairwise computation, so no numbers are invented.
% =====================================================================
\section{Two Workhorse Families}\label{sec:c30-apps}

\begin{previouslybox}[The machine is built -- now the two jobs it exists for]
Sections 26.1--26.2 gave the compressed tree and the weighted fold.  Almost every
virtual-tree problem in the wild is one of two shapes: a sum \emph{over all pairs}
of the marked set (distances, contributions), or a \emph{connected subgraph} spanning
the marked set (a Steiner tree, a count of ways to keep it connected).  Both collapse
to a single bottom-up pass over \(A(S)\), and both are worked here on the running
example so you can see the arithmetic before trusting it.
\end{previouslybox}

\begin{topicintro}
The reason these two families dominate is that a virtual tree preserves exactly two
things: the distance between any two marked nodes, and which marked nodes sit on which
side of every edge.  A pairwise sum decomposes over edges (each edge separates the
marked set into two groups, and \(\mathit{cnt}\cdot(k-\mathit{cnt})\) pairs cross it);
a connected spanning subgraph is forced to include every virtual edge, so its size is
a sum of the compressed lengths.  Neither needs the real tree at all.

\textbf{Prerequisites.} \(\secref{c30-construction}\), \(\secref{c30-dp}\),
\(\thmref{c30-fold}\).

\textbf{After this section you can}
\begin{itemize}[nosep]
  \item sum \(\dist\) over all pairs of a marked set in \(\bigO{k}\) after the build; \hfill\lv{5}
  \item read the size and cost of the minimal Steiner tree of a marked set off its virtual tree; \hfill\lv{5}
  \item count connected marked-induced structures with a small fold on \(A(S)\). \hfill\lv{5}
\end{itemize}
\end{topicintro}

% ---------------------------------------------------------------
\subsection{Family 1: the sum of pairwise distances}
% ---------------------------------------------------------------

\begin{theorem}{Pair-sum by per-edge separation}{c30-pairsum}
Let \(A(S)\) be the virtual tree of a marked set \(S\), \(|S|=k\), rooted at its top
LCA.  For a virtual edge \(e=(u,v)\) (with \(v\) the deeper end) let
\(\mathit{cnt}(v)\) be the number of marked nodes in \(v\)'s virtual subtree and
\(\mathit{w}(e)=\depth(v)-\depth(u)\) its real length.  Then
\[
  \sum_{\{x,y\}\subseteq S}\dist(x,y)
  \;=\; \sum_{e=(u,v)\in A(S)} \mathit{w}(e)\cdot \mathit{cnt}(v)\cdot\bigl(k-\mathit{cnt}(v)\bigr).
\]
\end{theorem}
\begin{proofbox}
In the real tree, edge \(e\) splits the vertices into the below-\(v\) side and the rest;
a pair \(\{x,y\}\) uses \(e\) on its path iff exactly one of \(x,y\) is below \(v\).  The
compressed edge represents \(\mathit{w}(e)\) identical real edges (a swallowed chain) with the
same below-side terminal count, so each of the \(\mathit{w}(e)\) edges is crossed by
\(\mathit{cnt}(v)(k-\mathit{cnt}(v))\) pairs; since \(\dist(x,y)\) is the number of edges on
the path, summing crossings over all edges counts every pair's distance exactly once.
The virtual tree's edges partition all marked-to-marked paths, so the sum over
\(A(S)\) equals the sum over the real tree. \(\Box\)
\end{proofbox}

\begin{tracebox}[Pair-sum on \(S=\{4,8,6,10\}\), \(k=4\): each virtual edge's share]
Using the virtual tree of \(\figref{c30-tree}\) with counts \(\mathit{cnt}(4){=}1\),
\(\mathit{cnt}(8){=}1\), \(\mathit{cnt}(2){=}2\), \(\mathit{cnt}(6){=}1\),
\(\mathit{cnt}(10){=}1\), \(\mathit{cnt}(3){=}2\):\\
\(2\!\to\!4\): \(1\cdot1\cdot3=3\). \quad \(2\!\to\!8\): weight\(2\), \(2\cdot1\cdot3=6\).\\
\(1\!\to\!2\): \(1\cdot2\cdot2=4\). \quad \(3\!\to\!6\): \(1\cdot1\cdot3=3\).\\
\(3\!\to\!10\): weight\(2\), \(2\cdot1\cdot3=6\). \quad \(1\!\to\!3\): \(1\cdot2\cdot2=4\).\\
Total \(=3+6+4+3+6+4=\boxed{26}\).  Direct check over the six pairs
(\(\dist(4,8)=3,\dist(4,6)=4,\dist(4,10)=5,\dist(8,6)=5,\dist(8,10)=6,\dist(6,10)=3\))
sums to \(3+4+5+5+6+3=26\). \(\Box\)
\end{tracebox}

\begin{lstlisting}[caption={Sum of pairwise distances over a marked set via its virtual tree. One bottom-up count, then the per-edge separation formula. ASCII only}, label={lst:c30-pairsum}]
// vtKids[v] = (child, length L); marked[v] = 1 if v is in S
ll cntv[MAXN];
void countDfs(int v){ cntv[v]=marked[v]; for(auto[u,L]:vtKids[v]) countDfs(u), cntv[v]+=cntv[u]; }
ll pairSum(int v, int k, ll acc){          // acc passed by reference in real code
    for (auto [u,L] : vtKids[v]){
        acc += (ll)L * cntv[u] * (k - cntv[u]);   // each of the L edges separates cntv[u] | k-cntv[u]
        pairSum(u,k,acc);
    }
    return acc;
}
// usage: buildVT(S); countDfs(rootOfVT); acc=0; ans=pairSum(rootOfVT,S.size(),acc);
\end{lstlisting}

% ---------------------------------------------------------------
\subsection{Family 2: the Steiner tree and connected structures}
% ---------------------------------------------------------------

On a tree the \term{Steiner tree} of \(S\) -- the minimum connected subgraph touching
all of \(S\) -- is unique: it is exactly the union of the paths between marked nodes,
i.e.\ the virtual tree with its compressed edges expanded.

\begin{itemize}[leftmargin=1.4em]
  \item \textbf{Size / edge count.} \(|\text{Steiner}| = \sum_{e\in A(S)} \mathit{w}(e)\).
        For \(S=\{4,8,6,10\}\) that is \(1+2+1+1+2+1=\boxed{8}\) real edges -- and indeed
        the spanning subtree is \(\{1,2,3,4,5,6,7,8,10\}\) minus the leaf \(9\), which has
        eight edges.
  \item \textbf{Vertex-weighted Steiner cost.} If edges/vertices carry costs, sum the
        cost of each virtual edge's chain, charging the endpoints once (a kept node shared
        by two virtual edges must not be double-billed -- the fold counts nodes through
        \(\mathit{cnt}\) and edges through \(\mathit{w}\)).
  \item \textbf{Connected-subgraph counts.} ``How many ways to choose a connected set
        containing all of \(S\)?'' becomes a two-state fold on \(A(S)\): at each kept
        node combine children with \(g_v=(\prod_{\text{child}}(1+g_{\text{child}}))\)
        minus the invalid all-excluded case, exactly Ch.~13's counting template run on the
        compressed tree.
\end{itemize}

\begin{insightbox}[Why Steiner-on-a-\emph{tree} is free but Steiner-on-a-graph is not]
Because a tree has a unique path between every pair, the minimal connected subgraph
spanning \(S\) is forced -- the Steiner tree is just that union of paths, computable in
\(O(k\log k)\) from the virtual tree with no search.  On a \emph{general graph} the
terminals may meet at newly-created branch nodes, so the problem is NP-hard and solved
by the Dreyfus--Wagner bitmask DP over \emph{terminal subsets} (\(O(3^k n)\)) -- a
different tool entirely.  That distinction is the heart of \(\secref{c30-boundary}\),
and the mixed set's CF 152E is precisely a graph/Steiner problem where a virtual tree
does \emph{not} apply.
\end{insightbox}

\begin{warningbox}[Double-billing the kept nodes]
A node of \(A(S)\) that is an LCA of two colours or two query terminals is shared by
several virtual edges; when a cost is per \emph{vertex} you must charge it once, not once
per incident edge.  The edge-count formulas (\(\thmref{c30-pairsum}\), the Steiner size)
are per-edge and safe; switch to a node fold the moment the statement prices vertices,
and start the fold from \(\mathit{cnt}\ge1\) nodes so pure-LCA nodes are still counted.
\end{warningbox}

\begin{hardeningbox}
\begin{itemize}[nosep,leftmargin=1.4em]
  \item \textbf{Root the virtual tree at the top LCA} (\(\lca\) of the first and last
        \(\tin\)-sorted marked node); the pair-sum and counts assume one root.
  \item \textbf{\(\mathit{cnt}\) counts terminals, not real subtree size.} A virtual
        child with no marked node under it never appears (it was pruned), so
        \(\mathit{cnt}(v)\ge1\) always; do not seed it with the real subtree size.
  \item \textbf{Modular counting} for connected-subgraph problems: keep every multiply
        under the modulus, and the \((1+g)\) product is over \emph{virtual} children.
\end{itemize}
\end{hardeningbox}

\begin{recallbox}[26.3 in one breath]
\emph{Two jobs, one fold.  Pairwise distance sums are \(\sum \mathit{w}(e)\cdot
\mathit{cnt}(v)(k-\mathit{cnt}(v))\) over virtual edges (\(\thmref{c30-pairsum}\)) -- the
running example gives \(26\), matching all six direct pairs.  The Steiner tree of a
marked set on a \textbf{tree} is just the expanded virtual tree, so its size is
\(\sum\mathit{w}(e)\) (eight here) and connected counts are a two-state Ch.~13 fold on
\(A(S)\).  Steiner on a \textbf{graph} is the other tool -- bitmask DP -- which is
exactly where the next-but-one section sends it.}
\selftest{
  \item Give the per-edge term for the pairwise distance sum. \emph{\(\mathit{w}(e)\cdot\mathit{cnt}(v)(k-\mathit{cnt}(v))\), \(\thmref{c30-pairsum}\).}
  \item What is the Steiner tree of a marked set on a tree? \emph{The union of paths between marked nodes -- the expanded virtual tree, forced and unique.}
  \item Why does Steiner on a graph need a different method? \emph{Terminals may meet at new branch nodes; NP-hard, bitmask Dreyfus--Wagner DP (\(\secref{c30-boundary}\)).}
}
\end{recallbox}
